\documentclass[]{article}

\usepackage{amsmath}
\usepackage{multirow}
\usepackage{natbib}
\usepackage{graphicx} 


\begin{document}

This study employs a rolling-window backtesting framework to compare two dynamic covariance matrix estimation strategies for Minimum Variance Portfolio construction. The analysis is performed on a dataset of daily returns, with all estimations and portfolio optimizations re-calculated at each time step.

\subsection{Data and dynamic volatility modeling}
The dataset consists of a 1,000-day panel of daily returns for ten large-capitalization U.S. equities, including assets known for high idiosyncratic volatility. To explicitly model the time-varying nature of asset risk, we first filter the raw returns using a Generalized Autoregressive Conditional Heteroskedasticity (GARCH) model.

For each asset $i$ and at each time step $t$, a GARCH(1,1) model is fitted to the most recent 60 days of historical returns. This procedure generates a one-step-ahead forecast of the conditional standard deviation, $\hat{\sigma}_{i,t+1}$. The raw returns are then standardized to produce a series of innovations, $z_{i,t}$, for each asset:
\begin{equation}
z_{i,t} = \frac{r_{i,t}}{\hat{\sigma}_{i,t}}
\end{equation}
These GARCH-filtered innovations, which are approximately homoskedastic, serve as the input for the two competing covariance estimation models described below. The final covariance matrix of returns is then recovered by re-scaling the estimated covariance matrix of innovations with the GARCH volatility forecasts.

\subsection{Covariance matrix estimation}
We construct and compare two distinct estimators for the covariance matrix of asset returns within each 60-day rolling window.

\subsubsection{Structural two-factor model}
The first estimator is a structural factor model that decomposes the covariance matrix of innovations, $\Sigma_{z,t}$, into systematic and idiosyncratic components:
\begin{equation}
\Sigma_{z,t} = B_t \Omega_t B_t^T + \Psi_t
\end{equation}
where $B_t$ is the $N \times 2$ matrix of factor loadings, $\Omega_t$ is the $2 \times 2$ covariance matrix of the factors, and $\Psi_t$ is the $N \times N$ diagonal matrix of idiosyncratic variances.

The two factors are defined as follows:
\begin{enumerate}
    \item \textbf{Market Factor:} The first principal component of the innovations of the ten assets in the estimation window, designed to capture the dominant source of systematic risk.
    \item \textbf{Sector Factor:} A long-short portfolio constructed from the innovations of a technology-focused subset of assets versus the remaining assets.
\end{enumerate}
To simplify the model, the Sector Factor is orthogonalized with respect to the Market Factor using the Gram-Schmidt process. This ensures that the factor covariance matrix $\Omega_t$ is diagonal. The factor loadings in $B_t$ are estimated via Ordinary Least Squares (OLS) regression of each asset's innovations onto the two orthogonal factors. The idiosyncratic variances in $\Psi_t$ are the variances of the residuals from these regressions.

The final covariance matrix of returns is then reconstructed by re-introducing the GARCH-forecasted volatilities:
\begin{equation}
\Sigma_{factor, t} = \text{diag}(\hat{\sigma}_t) (B_t \Omega_t B_t^T + \Psi_t) \text{diag}(\hat{\sigma}_t)
\end{equation}
where $\hat{\sigma}_t$ is the vector of conditional standard deviations for all assets at time $t$.

\subsubsection{Ledoit-Wolf shrinkage estimator}
The second estimator is the statistical shrinkage model proposed by Ledoit and Wolf. This method regularizes the sample covariance matrix of innovations, $S_t$, by pulling it towards a highly structured target matrix, $F_t$. The resulting shrinkage estimator, $\hat{\Sigma}_{z,t}$, is a weighted average of the two:
\begin{equation}
\hat{\Sigma}_{z,t} = (1-\delta_t)S_t + \delta_t F_t
\end{equation}
The target matrix, $F_t$, is a constant correlation matrix, where the diagonal elements are the sample variances of the innovations and all off-diagonal elements are derived from the average pairwise correlation observed in the estimation window. The shrinkage intensity, $\delta_t \in [0, 1]$, is not a free parameter but is calculated analytically at each time step to minimize the expected squared error between the estimated and true covariance matrices. The final covariance matrix of returns is then recovered in the same manner as the factor model:
\begin{equation}
\Sigma_{shrinkage, t} = \text{diag}(\hat{\sigma}_t) \hat{\Sigma}_{z,t} \text{diag}(\hat{\sigma}_t)
\end{equation}

\subsection{Portfolio construction and evaluation metrics}
Using the covariance matrices generated by each method, we construct and evaluate long-only Minimum Variance Portfolios (MVPs).

\subsubsection{Portfolio construction}
At the end of each day $t$, we solve the following quadratic optimization problem for each of the two estimated covariance matrices, $\Sigma_t$:
\begin{align}
\min_{w_t} \quad & w_t^T \Sigma_t w_t \\
\text{subject to} \quad & w_t^T \mathbf{1} = 1 \\
& w_{i,t} \ge 0 \quad \text{for } i=1, \dots, N
\end{align}
The resulting optimal weight vector, $w_t$, is then held for the subsequent day, $t+1$.

\subsubsection{Evaluation metrics}
The performance and stability of the two strategies are assessed using three key metrics:
\begin{enumerate}
    \item \textbf{Realized Portfolio Variance:} The primary measure of out-of-sample performance. The realized variance for day $t+1$ is calculated using the weights determined at time $t$ and the actual returns observed on day $t+1$: $\sigma^2_{p, t+1} = w_t^T r_{t+1} r_{t+1}^T w_t$. We compare the average realized variance over the entire backtest period.
    \item \textbf{Covariance Matrix Condition Number:} A diagnostic for numerical stability. The condition number is the ratio of the largest to the smallest eigenvalue of the covariance matrix, $\kappa(\Sigma) = \lambda_{max} / \lambda_{min}$. A high condition number indicates that the matrix is ill-conditioned, making its inverse (and thus the resulting portfolio weights) highly sensitive to small estimation errors.
    \item \textbf{Portfolio Turnover:} To quantify the magnitude of rebalancing required by each strategy, we compute the daily turnover as the sum of absolute changes in portfolio weights: $\text{Turnover}_t = \sum_{i=1}^{N} |w_{i,t} - w_{i,t-1}|$.
\end{enumerate}

\end{document}
                